\small
\begin{longtable}{p{3.0cm} p{2.0cm} p{10cm}}
    \caption{Campos del formato de mensaje Syslog definido en RFC 5424}
    \label{tab:rfc5424_fields} \\

    \toprule
    \textbf{Campo} & \textbf{Longitud} & \textbf{Descripción} \\
    \midrule
    \endfirsthead

    \multicolumn{3}{c}{\textit{Tabla \thetable. Continuación}} \\
    \toprule
    \textbf{Campo} & \textbf{Longitud} & \textbf{Descripción} \\
    \midrule
    \endhead

    \midrule
    \multicolumn{3}{r}{\textit{Continúa en la siguiente página}} \\
    \endfoot

    \bottomrule
    \endlastfoot

    \texttt{PRI} &
    1--3 caracteres &
    Contiene la prioridad del mensaje, calculada a partir de los valores de
    \textit{Facility} y \textit{Severity}. Se representa mediante un valor
    numérico encerrado entre los caracteres ``\texttt{<}'' y ``\texttt{>}''. \\

    \texttt{VERSION} &
    1--3 caracteres &
    Indica la versión del formato de mensaje Syslog utilizado. Para RFC 5424,
    el valor definido actualmente es \texttt{1}. \\

    \texttt{TIMESTAMP} &
    Variable &
    Indica la fecha y hora en la que se generó el evento. Utiliza un formato
    basado en ISO 8601 e incluye información de zona horaria y,
    opcionalmente, fracciones de segundo. \\

    \texttt{HOSTNAME} &
    Variable &
    Identifica el dispositivo o sistema que originó el mensaje. Puede
    corresponder al nombre de host, nombre de dominio o dirección IP del
    emisor, dependiendo de la información disponible. \\

    \texttt{APP-NAME} &
    1--48 caracteres &
    Identifica la aplicación o proceso que originó el mensaje. Permite asociar
    el evento con un componente específico del sistema emisor. \\

    \texttt{PROCID} &
    1--128 caracteres &
    Identifica el proceso que generó el mensaje. En sistemas donde sea
    aplicable, puede corresponder al identificador de proceso (PID). \\

    \texttt{MSGID} &
    1--32 caracteres &
    Identifica de forma única el tipo o clase de mensaje dentro del contexto
    de la aplicación que lo genera. Facilita la clasificación y correlación
    de eventos. \\

    \texttt{STRUCTURED-DATA} &
    Variable &
    Contiene información estructurada asociada al evento. Está formada por
    uno o más elementos de datos estructurados que pueden incluir parámetros
    con nombre y valor, facilitando el procesamiento automático de los
    mensajes. \\

    \texttt{MSG} &
    Variable &
    Contiene el contenido específico del evento. Puede incluir información
    adicional generada por la aplicación y, a diferencia de
    \textit{STRUCTURED-DATA}, generalmente se presenta como texto libre.

\end{longtable}